% GENERATED FILE. Edit canonical lessons, then run npm run generate:books.
% canonical-source-hash: fnv1a64:62b52b2dee0ab941
% canonical-lessons: LA-C468-tingo, LA-C468-murex, LA-C468-ostrum, LA-C468-pigmentum, LA-C468-penicillus

\chapter{To dye, Purple shellfish, Purple dye, Paint, Paintbrush}
\label{ch:la-to-dye-purple-shellfish-purple-dye-paint-paintbrush}
\hlchaptermodality{468}

\begin{chapteropening}
\textbf{By the end of this chapter:} \emph{I can talk about dyeing and painting: dyeing cloth, the purple shellfish, purple dye, paint and a paintbrush.}

Complete the last lesson of chapter 468: I can talk about dyeing and painting: dyeing cloth, the purple shellfish, purple dye, paint and a paintbrush.
\end{chapteropening}

\section[tingō, tingere]{tingō, tingere — to dye, to colour; to wet, to dip}

\label{lesson:LA-C468-tingo}



\begin{quote}
Before the new one: say the Latin for a tanner, then the Latin for a furrier.
\end{quote}

\subsection*{You'll want to know: tingō, tingere}
\textbf{tingō, tingere} — \textquotedblleft{}to dye, to colour; to wet, to dip\textquotedblright{}.

\begin{quote}
Servus lānam in aquā calidā tingit.
\end{quote}

\begin{quote}
\emph{The slave dyes the wool in hot water.}
\end{quote}

\emph{tingō, tingere, tīnxī, tīnctum}, also spelled \emph{tinguō}. English \emph{tinge}, \emph{tint} and \emph{tincture} come from it.

\begin{cousinweb}
\emph{tingō} is inherited from PIE \emph{teng-} 'to dip, to soak' (EtymDB), and Greek \emph{tengō} 'to wet' is a cousin. English \emph{tinge}, \emph{tint} and \emph{tincture} come from it.
\end{cousinweb}

\begin{grammarlens}[title={What you've built}]
One more word for telling, explaining and describing at B1.
\end{grammarlens}

\subsection*{Your turn}
\begin{itemize}
\raggedright
  \item \emph{Say it:} \emph{tingō, tingere}
  \item \emph{Say it:} \emph{tingō, tingere}, once more
  \item \emph{Say it:} a sentence of your own with \emph{tingō, tingere}
  \item \emph{Recall:} say the Latin for a tanner, then the Latin for a furrier, then say \emph{tingō, tingere} again
\end{itemize}

\begin{tcolorbox}[breakable,colback=teal!4,colframe=teal!35!black,title={Before you move on}]
What does tingō, tingere mean? (\textquotedblleft{}To dye, to colour; to wet, to dip\textquotedblright{}.) Say it once more.
\end{tcolorbox}
\section[mūrex, mūricis]{mūrex, mūricis — the purple shellfish, murex; purple dye}

\label{lesson:LA-C468-murex}



\begin{quote}
Before the new one: say the Latin for a furrier, then the Latin for to dye.
\end{quote}

\subsection*{You'll want to know: mūrex, mūricis}
\textbf{mūrex, mūricis} — \textquotedblleft{}the purple shellfish, murex; purple dye\textquotedblright{}.

\begin{quote}
Ex mūrice color purpureus fit.
\end{quote}

\begin{quote}
\emph{A purple colour is made from the murex shell.}
\end{quote}

Third declension, masculine: \emph{mūrex, mūricis m.} Thousands of these sea snails went into a little purple dye. English keeps \emph{murex}.

\begin{cousinweb}
\emph{mūrex} is usually traced to Greek \emph{myax} 'sea mussel' (EtymDB); Lewis \& Short give the sense but no source, so the path is hedged.
\end{cousinweb}

\begin{grammarlens}[title={What you've built}]
One more word for telling, explaining and describing at B1.
\end{grammarlens}

\subsection*{Your turn}
\begin{itemize}
\raggedright
  \item \emph{Say it:} \emph{mūrex, mūricis}
  \item \emph{Say it:} \emph{mūrex, mūricis}, once more
  \item \emph{Say it:} a sentence of your own with \emph{mūrex, mūricis}
  \item \emph{Recall:} say the Latin for a furrier, then the Latin for to dye, then say \emph{mūrex, mūricis} again
\end{itemize}

\begin{tcolorbox}[breakable,colback=teal!4,colframe=teal!35!black,title={Before you move on}]
What does mūrex, mūricis mean? (\textquotedblleft{}The purple shellfish, murex; purple dye\textquotedblright{}.) Say it once more.
\end{tcolorbox}
\section[ostrum, ostrī]{ostrum, ostrī — purple dye; purple cloth}

\label{lesson:LA-C468-ostrum}



\begin{quote}
Before the new one: say the Latin for to dye, then the Latin for the purple shellfish.
\end{quote}

\subsection*{You'll want to know: ostrum, ostrī}
\textbf{ostrum, ostrī} — \textquotedblleft{}purple dye; purple cloth\textquotedblright{}.

\begin{quote}
Dīves vestem ostrō tīnctam gerit.
\end{quote}

\begin{quote}
\emph{The rich man wears a garment dyed with purple.}
\end{quote}

Second declension, neuter: \emph{ostrum, ostrī n.} A word of poetry for the dye and for cloth coloured with it.

\begin{cousinweb}
\emph{ostrum} was borrowed from Greek \emph{ostreon} 'oyster, shellfish' (EtymDB; Lewis \& Short), since the dye came from a shellfish.
\end{cousinweb}

\begin{grammarlens}[title={What you've built}]
One more word for telling, explaining and describing at B1.
\end{grammarlens}

\subsection*{Your turn}
\begin{itemize}
\raggedright
  \item \emph{Say it:} \emph{ostrum, ostrī}
  \item \emph{Say it:} \emph{ostrum, ostrī}, once more
  \item \emph{Say it:} a sentence of your own with \emph{ostrum, ostrī}
  \item \emph{Recall:} say the Latin for to dye, then the Latin for the purple shellfish, then say \emph{ostrum, ostrī} again
\end{itemize}

\begin{tcolorbox}[breakable,colback=teal!4,colframe=teal!35!black,title={Before you move on}]
What does ostrum, ostrī mean? (\textquotedblleft{}Purple dye; purple cloth\textquotedblright{}.) Say it once more.
\end{tcolorbox}
\section[pigmentum, pigmentī]{pigmentum, pigmentī — a paint, a colour, a pigment}

\label{lesson:LA-C468-pigmentum}



\begin{quote}
Before the new one: say the Latin for the purple shellfish, then the Latin for purple dye.
\end{quote}

\subsection*{You'll want to know: pigmentum, pigmentī}
\textbf{pigmentum, pigmentī} — \textquotedblleft{}a paint, a colour, a pigment\textquotedblright{}.

\begin{quote}
Pictor pigmenta in tabulā miscet.
\end{quote}

\begin{quote}
\emph{The painter mixes his paints on a board.}
\end{quote}

Second declension, neuter: \emph{pigmentum, pigmentī n.}, from \emph{pingō} 'to paint'. English \emph{pigment} comes from it.

\begin{cousinweb}
\emph{pigmentum} is built on \emph{pingō} 'to paint' with the suffix \emph{-mentum} (EtymDB; Lewis \& Short). English \emph{pigment} comes from it.
\end{cousinweb}

\begin{grammarlens}[title={What you've built}]
One more word for telling, explaining and describing at B1.
\end{grammarlens}

\subsection*{Your turn}
\begin{itemize}
\raggedright
  \item \emph{Say it:} \emph{pigmentum, pigmentī}
  \item \emph{Say it:} \emph{pigmentum, pigmentī}, once more
  \item \emph{Say it:} a sentence of your own with \emph{pigmentum, pigmentī}
  \item \emph{Recall:} say the Latin for the purple shellfish, then the Latin for purple dye, then say \emph{pigmentum, pigmentī} again
\end{itemize}

\begin{tcolorbox}[breakable,colback=teal!4,colframe=teal!35!black,title={Before you move on}]
What does pigmentum, pigmentī mean? (\textquotedblleft{}A paint, a colour, a pigment\textquotedblright{}.) Say it once more.
\end{tcolorbox}
\section[pēnicillus, pēnicillī]{pēnicillus, pēnicillī — a paintbrush, a fine brush}

\label{lesson:LA-C468-penicillus}



\begin{quote}
Before the new one: say the Latin for purple dye, then the Latin for a paint.
\end{quote}

\subsection*{You'll want to know: pēnicillus, pēnicillī}
\textbf{pēnicillus, pēnicillī} — \textquotedblleft{}a paintbrush, a fine brush\textquotedblright{}.

\begin{quote}
Pictor pēnicillō tenuī flōrēs pingit.
\end{quote}

\begin{quote}
\emph{The painter paints flowers with a fine brush.}
\end{quote}

Second declension, masculine: \emph{pēnicillus, pēnicillī m.} (also \emph{pēnicillum n.}). English \emph{pencil} comes from it: the first pencils were fine brushes.

\begin{cousinweb}
\emph{pēnicillus} is a diminutive of \emph{pēniculus} 'brush' (Lewis \& Short), which first meant 'a little tail', since brushes were made of ox-tails and horse-tails. English \emph{pencil} comes from it through French.
\end{cousinweb}

\begin{grammarlens}[title={What you've built}]
One more word for telling, explaining and describing at B1.
\end{grammarlens}

\subsection*{Your turn}
\begin{itemize}
\raggedright
  \item \emph{Say it:} \emph{pēnicillus, pēnicillī}
  \item \emph{Say it:} \emph{pēnicillus, pēnicillī}, once more
  \item \emph{Say it:} a sentence of your own with \emph{pēnicillus, pēnicillī}
  \item \emph{Recall:} say the Latin for purple dye, then the Latin for a paint, then say \emph{pēnicillus, pēnicillī} again
\end{itemize}

\begin{tcolorbox}[breakable,colback=teal!4,colframe=teal!35!black,title={Before you move on}]
What does pēnicillus, pēnicillī mean? (\textquotedblleft{}A paintbrush, a fine brush\textquotedblright{}.) Say it once more.
\end{tcolorbox}
